\section{Deviations and Threats to Validity}
\label{sec:threats}

\paragraph{Code patches.} The released re-implementation required eight fixes to run: six
version-skew bugs on the training path (including rollout log-probabilities that \foldgrpo{}'s
importance ratios require but that the rollout server did not return) and two evaluation-path
fixes. All are committed with documentation; training hyperparameters are the authors' script
verbatim.

\paragraph{Effective steps.} \grpo{} trained 98 effective steps of 100
(\cref{sec:results}); \foldgrpo{} trained 100.

\paragraph{Synchronous training.} The paper trains with asynchronous rollout, allowing up to
five steps of off-policy staleness between generation and update. The released
re-implementation does not include that pipeline: both of our arms train on-policy with a full
per-step barrier---generation, reward, advantage, and update run sequentially, and every step
waits for all 256 trajectories (including judge calls) to complete. Trajectory execution
\emph{within} the generation phase is concurrent, so tool and judge latency overlaps
generation across trajectories. Both arms share this property, so the comparison is
unaffected, but per-step wall-clock time is dominated by straggler trajectories and is not
comparable to the paper's throughput.

\paragraph{Judge substitution.} \judgemodel{} replaces the paper's proprietary judges in all
three roles, so absolute numbers are not comparable to the paper's; within-study comparisons
use one fixed judge, audited manually on 49 calls.

\paragraph{Statistical power.} One seed per arm; the greedy tier has $n{=}150$ (CI
$\approx \pm 4$ points); one base-model scale (8B). The paper's 36B margin may be real at scale;
our result shows only that it does not manifest at 8B under a matched protocol.

\paragraph{Scope.} We replicate the \bcplus{} (deep research) track only; the paper's software
engineering track is not reproducible from the released code, which omits the corresponding
training environment.
